\clearpage
\section*{September 9, 2026: A second ten-case construction review}
\addcontentsline{toc}{section}{September 9, 2026: A second ten-case construction review}
\textit{Agent-drafted findings; recommendations not yet applied.}

Eight of the ten reviewed entries are old parent records with individually recorded homes already retained. These comprise two Stave groups, four Enclave Bucktown groups, Hermitage and Janssen. The recommendation is to retire the parents and preserve the separate homes, rather than add the parents as more construction. The eighteen Stave homes are already represented by eighteen projects, including nine homes that each occupy two tax parcels. The four Enclave groups link to 46 homes; two other retained houses bring the development to 48. Original permits describe 49 homes, but completed revision 100699119 reduces one building from seven to six. The retained 48 therefore agree with the revised permits. No home was created from a marketing total.

The Sheffield permits establish two separate buildings with three apartments each and rooftop club space. The Assessor's four-unit entries should not become four apartments per building. The operator reports construction in 2008, versus the Assessor's 2006. These count and year corrections are recommendations pending adoption; floor-area coverage still needs confirmation.

At Ashland, three completed permits and the marketing account identify eighteen apartments in three buildings. A fourth Assessor card on the second parcel should not be assumed to be another building. Marketing reports completion in 2020, versus the Assessor's 2021. The common site's land coverage and repeated card still need a final measurement decision. At Hermitage, permits issued in October 2016 rule out the old parent's 2014 year; the individually retained years are 2016 and 2017. Janssen's current listing includes an additional side lot, so its marketed lot area must not replace the individual Assessor denominator.

\medskip
\noindent\textit{Evidence and limits.} The release-review notes record all ten source IDs, retained counterparts and public links. Findings use the pinned Assessor and full permit histories. The only source record in this batch in the frozen 500-foot sample is Janssen. No production data or estimate changed; the fixed review list remains 196 closed and 72 open until decisions are applied and verified.

\clearpage
\subsection*{September 9 follow-up: Sheffield and Ashland measurements}
The Sheffield unit discrepancy is explainable: the official Assessor page labels four \emph{total} units, and the saved record identifies one commercial unit. This agrees with each new-building permit's three apartments and rooftop club. Each parcel has its own reported 7,308-square-foot mixed-use building area; the inspected official page has no separate commercial-building area to add. The recommendation is two buildings, each with three homes, completed in 2008 according to the operator, using the reported building areas and individual Assessor land areas of 4,968 square feet at 3617 and 3,000 at 3619. These are source measurements, not independently reconstructed residential-only areas.

At Ashland, the current broker listing identifies exactly the two Assessor parcels and eighteen apartments completed in 2020. Its 16,988-square-foot site differs by only 87 square feet from the Assessor parcels' combined 16,901. Three six-apartment building permits support counting the primary parcel's three building cards once; the second parcel repeats a card under the same primary property reference. The three areas total 45,543 square feet. An older advertisement reports 38,970, but also gives a different parcel number, a 2000 construction year and a much smaller lot. That advertisement should not override the Assessor measurements.

The recommendation is one eighteen-home Ashland project containing three buildings, with 2020 completion, 16,901 square feet of land and 45,543 square feet of building area. Separate building lots are not reported, so no land allocation is proposed. The floor-area recommendation rests on the Assessor cards and three permits, not independent confirmation of the total. Jacob favors retaining 2016 for 5933 Hermitage; that is already its selected Assessor year. Sheffield and Ashland recommendations have not been applied, and no dataset or estimate changed in this follow-up. Public source links and parcel-level evidence are recorded in the release-review notes on the working branch.
